\scanpage{115}
\mspage{\handwritten{85}}

sequence

\diagram{\handwritten{\begin{tikzcd}[ampersand replacement=\&,column sep=small]
0\arrow[r] \& O_V(-\nu  E')\arrow[r] \& O_V(-\mathscr E'_\nu)\arrow[r] \& G\arrow[r] \& 0\unsure{,|.}
\end{tikzcd}}}

Here $G$ is some coherent sheaf, concentrated on $\handwritten{\mathscr E'}_\nu$.
It follows from the observation made in the beginning of
the proof that, beginning with some sufficiently large $n$,

\[
\handwritten{H^1(V,G[nE'])=0.}
\]

Hence it follows that, beginning with some sufficiently
large $n$,

\[
\handwritten{\dim H^1(V,O_V(-\mathscr E'_\nu)[nE'])\leq\dim H^1(V,O_V[(n-\nu)E']).}\tag*{\handwritten{(21)}}
\]

or

\[
\handwritten{\dim H^1(V,O_V(-\mathscr E'_\nu)[nE'])\leq\dim H^1(V,O_V[(n-\nu)E']).}
\]

From the exact sequence

\diagram{\handwritten{\begin{tikzcd}[ampersand replacement=\&,column sep=small]
0\arrow[r] \& O_V(-\mathscr E'_\nu)\arrow[r] \& O_V\arrow[r] \& O_{\mathscr E'}\arrow[r] \& 0
\end{tikzcd}}}

we obtain, at $n\geq n(\handwritten{\mathscr E'}_\nu)$:

\[
\handwritten{\dim H^1(V,O_V[nE'])\unsure{\leq|\leqq}\dim H^1(V,O_V(-\mathscr E'_\nu)[nE']),}\tag*{\handwritten{(22)}}
\]

% archive sheet number: 114
